

\begin{essay}{}
    Ключевые слова:
    3Д---реконструкция,
    построение карты глубин,
    методы структурного освещения (Shape From Shading),
    определения глубины из размытия (Depth From Defocus),
    комфортной генерации глубины (Comfortable Depth Generation),
    на основе нейронных сетей (MiDaS).

    Объект исследования --- восстановление трехмерной информации из монокулярных изображений.

    Предмет исследования --- метод восстановления глубины из монокулярного изображения,
    основанный на комбинации нейросетевого подхода и классического метода восстановления глубины по закраске изображения.


     Целью работы является разработка гибридного метода восстановления
     трехмерной информации из архивных фотографий.

    Для достижения поставленной цели решены следующие задачи:
    \begin{itemize}
        \item проанализированы существующие методы восстановления глубины из
        монокулярных изображений;
        \item разработана формальная модель гибридного метода, объединяющего нейросетевой подход и классический метод определения глубины по закраске изображения;
        \item реализован программный прототип предложенного метода с использованием современных технологий обработки изображений;
        \item проведено исследование влияния степени состаривания изображения на точность полученной карты глубин;
        \item сформулированы рекомендации по применению предложенного метода.
    \end{itemize}


    В ходе исследования установлено дифференцированное влияние различных типов деградации изображений на точность восстановления глубины:
     структурные изменения (шум, виньетирование) оказывают наибольшее негативное влияние, а цветовая деградация (сепия, выцветание) влияет минимально.

    Область применения разработанного метода --- восстановление трёхмерной информации из архивных фотографий для выполнения задач цифровой реставрации и трёхмерной реконструкции исторических объектов
     при наличии фотоматериалов различного качества с минимальным уровнем шумовых искажений.
    В таких условиях гибридный метод обеспечивает баланс между устойчивостью
    к деградации и детализацией, что подтверждено проведенным исследованием:
    точность реконструкции нейросетевого алгоритма повышается на 1--2\% без необходимости переобучения модели.

    Перспективными направлениями дальнейшего развития разработанного метода являются:
    \begin{itemize}
        \item расширение функциональности для построения полных трехмерных моделей на основе одиночных архивных фотографий,
        \item интеграция с системами дополненной и виртуальной реальности для создания интерактивных исторических реконструкций,
        \item разработка специализированных алгоритмов предварительной реставрации изображений, оптимизированных для последующего восстановления глубины.
    \end{itemize}



%
%    Данная научно--исследовательская работа включает
%    анализ предметной области,
%    описание основных подходов к решению задачи конвертации 2Д изображений в 3Д,
%    определение критериев сравнения,
%    и классификацию методов.
\end{essay}